\section{Related work}
\label{sec:related-work}

% WRITING PROMPTS (comments only; they do not print). About 1 page, three
% subsections or three paragraph groups. For each group, end with what the
% group does NOT do that your paper does.
%
% Coverage path planning. Start from the surveys (Galceran and Carreras 2013;
% Liu et al. 2024 for mower robots). Boustrophedon cell decomposition (Choset
% 2000), the minimum-width sweep direction (Huang 2001), Fields2Cover.
% Question to answer: who chooses the sweep direction, and by what criterion?
% (usually the number of lanes or turns, not energy).
%
% Soft computing for ordering sub-regions. GA, PSO and ACO formulations as a
% travelling salesman problem over cells. Present Xie et al. (2024) in detail,
% because it is the baseline: GA then ant colony, centre-point tour, length and
% turns measured after planning. Which papers decode entry corners exactly?
%
% Energy-aware planning for skid-steer robots. The skid-steer power models
% (Wong; Yu et al. 2009; Morales et al.). Energy-aware coverage for UAVs and
% field robots. Point out that few CCPP papers put a skid-steer turning model
% inside the planning objective.
%
% Fuzzy path tracking (short): fuzzy controllers for mobile robot path
% following, to motivate the tracker of Fig. fuzzy-controller.
%
% Close with a table or a paragraph that places this paper against the
% closest three or four works on: objective, direction choice, robot model,
% real robot.
